% =====================================================================
% CSE 341: Concepts of Programming Languages
% Project D5 (Part 1): Retrospective & Self-Assessment for ANS
% Student: 220104004931 Enes Özkan
% Compile with: tectonic 220104004931_D5_P1.tex
% =====================================================================
\documentclass[11pt,a4paper]{article}

\usepackage[utf8]{inputenc}
\usepackage[T1]{fontenc}
\usepackage{mathpazo}
\usepackage[scaled=0.92]{helvet}
\usepackage[varqu,varl]{zi4}
\usepackage[english]{babel}
\usepackage{microtype}
\usepackage[a4paper,margin=2.25cm,top=2.0cm,bottom=2.15cm]{geometry}

\usepackage[table]{xcolor}
\definecolor{ansblue}{HTML}{1E3A8A}
\definecolor{anscaption}{HTML}{6B7280}
\definecolor{anscodebg}{HTML}{F9FAFB}
\usepackage[colorlinks=true, linkcolor=ansblue, urlcolor=ansblue,
            citecolor=ansblue, pdfborder={0 0 0}]{hyperref}

\usepackage[explicit]{titlesec}
\titleformat{\section}
  {\large\bfseries\sffamily\color{black}}
  {\thesection}{0.65em}{#1}
\titlespacing*{\section}{0pt}{11pt}{5pt}
\setlength{\parskip}{4.5pt plus 1pt minus 1pt}
\setlength{\parindent}{0pt}
\linespread{1.08}
\setcounter{secnumdepth}{0}

\usepackage{fancyhdr}
\pagestyle{fancy}
\fancyhf{}
\renewcommand{\headrulewidth}{0pt}
\fancyfoot[L]{\footnotesize\color{anscaption}%
  CSE 341: ANS Retrospective (D5, Part 1) \(\cdot\)
  Student 220104004931 Enes Özkan}
\fancyfoot[R]{\footnotesize\color{anscaption}Page \thepage}

\newcommand{\code}[1]{\texttt{#1}}

\hypersetup{
  pdftitle  = {ANS Retrospective and Self-Assessment (D5, Part 1)},
  pdfauthor = {220104004931 Enes Özkan},
  pdfsubject= {CSE 341 Project Part 1},
  pdfkeywords={DSL, chemistry, retrospective, self-assessment, ANS}
}

\begin{document}

\begin{center}
  {\footnotesize\color{anscaption}CSE 341: Concepts of Programming Languages\par}
  \vspace{3pt}
  {\huge\bfseries\sffamily ANS: Alchemist's Notation Script\par}
  \vspace{4pt}
  {\large Retrospective \& Self-Assessment (D5), Part 1\par}
  \vspace{6pt}
  {\footnotesize\color{anscaption}
    Student: 220104004931 Enes Özkan \(\cdot\)
    Spring 2026 \(\cdot\)
    Gebze Technical University\par}
\end{center}
\vspace{7pt}

\section{What Is Working}

At this point I trust the foundation more than I did at the start, but I do not want to oversell it. The C
lexer and recursive-descent parser pass the current valid and malformed D3 tests, and the diagnostics include
file position, which makes the errors usable. The AST dump is also helpful because it lets me compare the
implementation against the EBNF instead of guessing from the source code alone. The part I like most is that
the chemistry notation still reads naturally: \code{compound water}, \code{H2O}, and \code{CH4 + 2*O2}
all keep the program close to the problem domain.

\section{What Is Hard}

The hard part is not only writing the grammar, but making the lexer and parser behave exactly like that
grammar. Formula literals and identifiers look similar enough that a small lexical decision can affect later
syntax. Units are another example. \code{890.3 kJ/mol} is natural chemistry notation, but the same slash is
also an arithmetic operator. Implementing this in C made the work more explicit: AST nodes, recovery paths,
and memory ownership all have to be handled directly. That is useful for learning, but it also makes small
design changes expensive.

\section{Plan For Part 2}

For Part 2, the priority is execution without drifting away from the design. I need symbol tables, static
scoping, type checking, and an interpreter that respects the precedence and assignment rules from D1. One
specific design fix is output: Part 1 treats \code{print} as a special statement, but it would fit ANS better as
a normal built-in call, in the same style as other helper functions. If we make that change, D1, D2, and D3
must all be updated together. I also need domain-specific checks for reaction and unit mistakes, real outputs
for the D3 programs, and one type error caught before execution.

\end{document}
